\section{Detailed Hyperparameter Settings}
\label{appendix:hyperparameters}

In this section, we provide a detailed overview of the hyperparameter settings used across all experiments. All experiments were conducted on NVIDIA A800 GPUs.

\noindent\textbf{Reasoning and Coding Tasks.} For tasks involving logical reasoning (GSM8K, MATH500) and code generation (HumanEval, MBPP), we set the token sampling temperature $\tau_{\text{token}} = 0.7$ to strike a balance between output diversity and structural coherence. For the Info-Gain Sampler, we employ a relatively low position temperature $\tau_{\text{pos}} = 0.1$ to focus the selection on the most informative candidate positions, while generating $N=8$ candidate actions per step for parallel evaluation. To maintain high throughput during predictable decoding phases, we set the acceleration threshold $\gamma = 0.8$, which allows the sampler to skip the full evaluation routine when the maximum token probability is high. For block diffusion settings on Semi-AR models (SDAR and TraDo), we use a fixed block size of 16. The decoding budget $K$ (tokens per step) is varied between 1 and 2 to evaluate performance under different acceleration ratios. Maximum generation lengths are benchmark-specific: 256 tokens for GSM8K, HumanEval, and MBPP; 512 tokens for MATH500 and SDAR benchmarks; and 1024 tokens for the TraDo-8B model to accommodate longer reasoning chains.

\noindent\textbf{Text-to-Image Generation.} For multimodal experiments using the MMaDa model on ImageNet and GenEval, we adopt a more conservative token temperature $\tau_{\text{token}} = 0.4$ and a matching position temperature $\tau_{\text{pos}} = 0.4$ to ensure high fidelity and alignment with text prompts. The candidate set size is kept at $N=8$. We follow a 50-step decoding trajectory governed by a cosine schedule, which progressively reduces the number of masked positions to refine image details. For extreme acceleration tests (Section~\ref{appendix:few_step}), we utilize a 5-step linear schedule to evaluate the sampler's robustness under severe budget constraints.

\noindent\textbf{Creative Writing.} We evaluate creative writing performance using 200 prompts selected from the Alpaca dataset. Following the Min-P paper~\cite{nguyen2024turning}, we employ GPT-4.1~\cite{openai_gpt4_1} as an LLM judge to assess generation quality. For this experiment, the maximum generation length is set to 1024 tokens, with a fixed block size of 16.

% \noindent\textbf{Baselines} All baseline samplers (Uniform, Confidence, Entropy, Margin, KLASS, and PC-Sampler) are configured with the same token temperature $\tau_{\text{token}}$ as our method to ensure a fair comparison. For KLASS, we set the consistency threshold $\epsilon = 5 \times 10^{-4}$ following its original implementation. For PC-Sampler, we adopt the default decay parameters for positional weighting as specified in the original work.
